\documentclass[10pt,conference]{IEEEtran}
\IEEEoverridecommandlockouts

\usepackage{cite}
\usepackage{amsmath,amssymb}
\usepackage{algorithmic}
\usepackage{graphicx}
\usepackage{textcomp}
\usepackage{xcolor}
\usepackage{booktabs}
\usepackage{url}
\usepackage{hyperref}
\hypersetup{
    colorlinks=true,
    linkcolor=black,
    citecolor=black,
    urlcolor=blue,
}

\def\BibTeX{{\rm B\kern-.05em{\sc i\kern-.025em b}\kern-.08em
    T\kern-.1667em\lower.7ex\hbox{E}\kern-.125emX}}

\begin{document}

\title{Cross-Geo Product Matching with a\\Multimodal Siamese Network and CLIP-KNN Clustering}

\author{\IEEEauthorblockN{Matouš Kovář}
\IEEEauthorblockA{\textit{FIT, Czech Technical University in Prague} \\
kovar@viagem.cz}}

\maketitle

\begin{abstract}
We address the Glami 2026 cross-geo product matching challenge: identifying
identical products across thirteen European marketplaces despite differing
titles, descriptions, prices and images. We combine a frozen vision backbone
(CLIP-ViT-B/32) with a frozen multilingual text encoder
(CLIP-ViT-B-32-multilingual-v1) and structured metadata (price, geo,
department, color). A small Siamese head trained with binary cross-entropy on
hard-mined pairs scores each candidate; clustering is performed by cosine
$k$-nearest-neighbour retrieval followed by Siamese re-scoring and connected
components. On a 30k held-out slice of \texttt{items\_train} we achieve a
pairwise clustering $F_1 \approx 0.82$ (Phase~2 metric); on the Phase~1
binary task the same model attains $F_1 \approx 0.97$. We compare against
exact-title and raw-CLIP-cosine baselines, and discuss the failure mode that
motivated a (currently) unfinished metric-learning extension.
\end{abstract}

% ============================================================================
\section{Introduction}
\label{sec:intro}

% TODO: Tighten or expand to one column max.
Cross-geo product matching is the task of partitioning a catalogue of items
listed across multiple marketplaces into groups, where each group represents
a single underlying product (SKU). It is a special case of multimodal entity
resolution. The setting in this work is the Glami 2026 challenge: roughly
$1.3$~M items spanning thirteen European geographies (Czech, Slovak,
Hungarian, Romanian, Bulgarian, Croatian, Greek, Italian, Polish, Latvian,
Lithuanian, Slovenian, Estonian) with $217{,}000$ distinct product labels in
the training set.

Three properties make the problem difficult:

\textbf{Multimodality.} Listings include short, often noisy titles and longer
descriptions in the local language, a single product image (sometimes a
generic stock photo shared across SKUs), a numeric price in the local
currency, a coarse \texttt{departmentIds} category, and a multi-valued
\texttt{colorTagIds} field. No single modality is sufficient---two listings
may share an image but differ in price/SKU, or share a title but show a
different product.

\textbf{Cross-geo variation.} The same product is described in different
languages, photographed differently, and priced in different currencies.
Naive within-language matching fails by construction.

\textbf{Class imbalance.} Most products in
\texttt{items\_train} have only one or two listings. Among $217{,}000$
labels, roughly $54\%$ are singletons. Sampling random pairs as negatives is
trivially solvable; the discriminative signal lives in pairs that look
similar but are not duplicates.

The challenge has two phases, with two metrics:
(i) Phase~1 is binary classification over fifteen-thousand groups of five
items, scored by group-level $F_1$;
(ii) Phase~2 is end-to-end clustering on $\approx 200{,}000$
unlabeled items, scored by \emph{pairwise} $F_1$ over all C($n$,2) item pairs.

\textbf{Contribution.} We present a simple late-fusion Siamese architecture
that augments a learned $|e_1 - e_2|$ embedding-difference vector with raw
modality-level similarity features (image cosine, text cosine, Jaccard
overlap of department/color, geo match). We train it with embedding-mined
hard positives and hard negatives, and we re-use the same scorer for both
phases via a thin clustering wrapper based on cosine $k$-NN retrieval
followed by connected components. Despite its simplicity, the system attains
Phase-1 $F_1 \approx 0.97$ and a Phase-2 pairwise $F_1$ comparable to the
strongest baselines we evaluated.

% ============================================================================
\section{Related Work}
\label{sec:related}

% TODO: Confirm Zalando paper citation and add 2-3 more.
The closest prior work is the Zalando cross-geo matching study
\cite{zalando_crossgeo}, which targets the same problem on Zalando catalogue
data. Their pipeline uses a fashion-domain-specific image encoder, language
features, and a deep matching head; reported $F_1$ is in the high-$0.9$s for
top-$k$ retrieval but they report no end-to-end clustering pairwise $F_1$
that is directly comparable to our Phase~2 metric.

Our work differs in three respects:
(1) we keep the encoders frozen and train only a small fusion + classifier
head, which trains in minutes on consumer hardware;
(2) we explicitly inject \emph{raw} modality-level similarity features into
the classifier, which (as we discuss in \S\ref{sec:results}) is a
double-edged sword;
(3) we use multilingual CLIP text embeddings rather than fine-tuning a
language model, which sidesteps the need for in-domain text supervision.

CLIP \cite{clip} provides our vision and text backbones. FashionCLIP
\cite{fashionclip} is a domain-tuned variant that we briefly evaluated for
the metric-learning extension (\S\ref{sec:results}). Triplet/InfoNCE-style
metric learning \cite{facenet,supcon} is an obvious alternative to pair
classification; we leave a full evaluation to future work.

Classical entity-resolution baselines \cite{er_survey} (exact-match
hashing, TF--IDF~+ cosine, blocking) remain competitive on text-heavy
catalogues; we use them as our reference baselines (\S\ref{sec:baselines}).
Our clustering design---build a $k$-nearest-neighbour graph over learned
embeddings, edge-weight or filter via a learned scorer, then extract
connected components---is the standard pattern for large-scale face/identity
deduplication. Otto, Wang and Jain \cite{otto_facecluster} cluster millions
of faces by constructing a KNN graph over FaceNet embeddings and running
approximate Rank-Order clustering on it; Yang \emph{et al.}\
\cite{yang_lcluster} use the same KNN-graph substrate but learn a graph
convolutional network to predict cluster membership. We follow the same
two-stage recipe (KNN retrieval, then graph-level clustering) but keep the
clustering step deliberately simple---a Siamese score threshold followed by
connected components---because pair-supervision is what we have. For the
retrieval step itself, an approximate index such as FAISS \cite{faiss} would
be the production choice; we use exact block-matmul KNN since $N \le 200{,}000$
fits in memory.

% ============================================================================
\section{Methodology}
\label{sec:methodology}

\subsection{Features}
For every item $x$ we precompute and cache:
\begin{itemize}
\item $\mathbf{e}_x^{\text{img}} \in \mathbb{R}^{768}$, the
\texttt{pooler\_output} of \texttt{openai/clip-vit-base-patch32};
\item $\mathbf{e}_x^{\text{txt}} \in \mathbb{R}^{512}$, the output of
\texttt{sentence-transformers/clip-ViT-B-32-multilingual-v1} on the
concatenation of \texttt{title} and \texttt{description};
\item $p_x \in \mathbb{R}$, the price after a per-geo \texttt{StandardScaler}
fit on \texttt{items\_train};
\item one-hot/multi-hot indicators for $\texttt{geo}$, $\texttt{departmentIds}$
and $\texttt{colorTagIds}$.
\end{itemize}

\subsection{Embedder}
The shared embedder projects modality features into a joint space:
\[
\mathbf{z}_x = \texttt{MLP}\!\left(
  \texttt{LN}(W_i \mathbf{e}_x^{\text{img}}) \,\Vert\,
  \texttt{LN}(W_t \mathbf{e}_x^{\text{txt}}) \,\Vert\,
  p_x
\right) \in \mathbb{R}^{128},
\]
with $W_i \in \mathbb{R}^{256\times 768}$, $W_t \in \mathbb{R}^{256\times
512}$ and dropout $0.2$/$0.3$ between layers.

\subsection{Pair classifier}
We compute a pair feature vector that combines a learned interaction with
\emph{raw} similarity readouts:
\[
\mathbf{f}_{xy} = \big[\,
  |\mathbf{z}_x - \mathbf{z}_y|;\;
  \langle \mathbf{e}_x^{\text{img}}, \mathbf{e}_y^{\text{img}}\rangle;\;
  \langle \mathbf{e}_x^{\text{txt}}, \mathbf{e}_y^{\text{txt}}\rangle;\;
  J_{\text{dept}};\;
  J_{\text{color}};\;
  \mathbb{1}[\text{geo}_x{=}\text{geo}_y]\,\big]
\]
where $J$ is Jaccard overlap of multi-hot vectors. A two-layer MLP maps
$\mathbf{f}_{xy} \in \mathbb{R}^{133}$ to a logit. Total trainable
parameters: $\approx 519$k. Loss: \texttt{BCEWithLogitsLoss}; optimiser:
Adam, $\eta=10^{-3}$, weight decay $10^{-4}$;
\texttt{ReduceLROnPlateau} on validation $F_1$.

\subsection{Pair mining}
Random negatives saturate the loss within one epoch. We instead mine pairs
using the joint similarity $s(x,y) = \tfrac{1}{2}(\langle
\hat{\mathbf{e}}_x^{\text{img}}, \hat{\mathbf{e}}_y^{\text{img}}\rangle +
\langle \hat{\mathbf{e}}_x^{\text{txt}},
\hat{\mathbf{e}}_y^{\text{txt}}\rangle)$:
\begin{itemize}
\item \textbf{Hard negatives:} pairs with $s(x,y) \ge 0.9$ and
$\text{label}_x \ne \text{label}_y$. With $1\%$ probability we also accept
$s < 0.9$ to avoid distribution collapse.
\item \textbf{Hard positives:} same-label pairs with $s(x,y) < 0.85$
(low-similarity true matches across geos/photos), again with a small
admission probability for easier pairs.
\end{itemize}
The split is \emph{label-disjoint}: $20\%$ of labels are reserved for
validation and never appear in any training pair.

\subsection{Phase~2 clustering}
At inference (Algorithm~\ref{alg:cluster}) we (i) build joint $1280$-d
$L_2$-normalised vectors $[\mathbf{e}^{\text{img}} \,\Vert\,
\mathbf{e}^{\text{txt}}]$ for every item, (ii) compute top-$K{=}40$ cosine
nearest neighbours in $512$-row blocks---following the
embedding-KNN-graph construction used in face clustering
\cite{otto_facecluster,yang_lcluster} but implemented with exact block-matmul
on MPS rather than an approximate index, since $N \le 200{,}000$ fits in
memory---and any neighbour with cosine $\ge 0.60$ is kept as a candidate edge, (iii) re-score
every candidate edge with the trained Siamese classifier and discard edges
below a tuned threshold $\theta^\ast$, and (iv) run
\texttt{scipy.sparse.csgraph.connected\_components} on the surviving edges.
Components are split into chunks of at most $100$ items per submission line.

\begin{algorithm}[H]
\caption{End-to-end clustering at inference.}
\label{alg:cluster}
\begin{algorithmic}[1]
\STATE Build $\mathbf{M} \in \mathbb{R}^{N \times 1280}$, $L_2$-normalised.
\STATE $E \leftarrow \emptyset$
\FOR{block $B$ of $\mathbf{M}$}
  \STATE $S \leftarrow B \mathbf{M}^\top$ \quad \emph{(cosine since normalised)}
  \STATE $E \leftarrow E \cup \{(i,j) : S_{ij} \ge 0.60,\, j \in \text{top-}40,\, j>i\}$
\ENDFOR
\STATE $E^\ast \leftarrow \{(i,j) \in E : \texttt{Siamese}(i,j) \ge \theta^\ast\}$
\STATE Return connected components of $G=(V,E^\ast)$ chunked at $100$.
\end{algorithmic}
\end{algorithm}

% ============================================================================
\section{Baselines}
\label{sec:baselines}

We compare against three increasingly informed reference points:

\textbf{B1 -- Singleton (no merging).} Every item is its own cluster.
Provides the pairwise $F_1$ baseline of $0$ (precision undefined by
convention); useful only as a sanity floor.

\textbf{B2 -- Exact-title match.} Merge items whose lowercased,
whitespace-stripped \texttt{title} string is identical. Strong on within-geo
duplicates but degrades sharply across languages.

\textbf{B3 -- Raw CLIP cosine.} Cluster directly on
$[\mathbf{e}^{\text{img}}\,\Vert\,\mathbf{e}^{\text{txt}}]$ via the same
KNN+connected-components pipeline as our method, but with \emph{no} Siamese
re-scoring---only the cosine threshold $\theta_\text{cos}$ is tuned. This
isolates the value added by the learned head.

% ============================================================================
\section{Experiments}
\label{sec:experiments}

\subsection{Validation strategy}
We hold out $20\%$ of \emph{labels} (not items) from \texttt{items\_train}
and use them to construct two evaluation sets:
\begin{itemize}
\item \textbf{Pair-F1 set.} 400k mined pairs (75\% hard negatives, 25\% hard
positives) over the held-out labels. Used for early stopping during training.
\item \textbf{Clustering set.} A 30k-item slice (15k items from held-out
labels, 15k random distractors). The full pipeline runs on this slice and we
report \emph{pairwise} $F_1$, the same metric used for Phase~2 grading.
\end{itemize}

\subsection{Hyperparameter selection}
The Siamese threshold $\theta^\ast$ is selected by sweep over the
clustering set. The \texttt{KNN\_COSINE\_THRESHOLD} ($0.60$) and
\texttt{TOP\_K} ($40$) were chosen so that the candidate-edge set has
\emph{high} pair recall (over $95\%$ of true positive pairs reachable) at
acceptable size ($\sim 600$k edges for $30$k items).
Table~\ref{tab:sweep} shows the full sweep over $\theta$. Pairwise
$F_1$ peaks sharply at $\theta^\ast = 0.97$ ($F_1 = 0.824$); the size of the
largest predicted cluster is the cleanest leading indicator of
over-merging---it stays under $50$ items down to $\theta = 0.95$ but jumps
past $200$ at $\theta = 0.90$ and past $500$ at $\theta = 0.85$, at which
point a few false-positive edges chain entire connected components together
and precision collapses while recall barely moves.

\begin{table}[t]
\centering
\caption{Siamese threshold sweep on the held-out 30k clustering slice. ``Edges''
is the number of pair candidates surviving threshold $\theta$; ``Max cl.''
is the size of the largest predicted connected component.}
\label{tab:sweep}
\begin{tabular}{lrrrrr}
\toprule
$\theta$ & Edges & $P$ & $R$ & $F_1$ & Max cl. \\
\midrule
0.99 & 13{,}691 & 0.978 & 0.674 & 0.798 & 18  \\
0.98 & 15{,}210 & 0.963 & 0.716 & 0.821 & 19  \\
\textbf{0.97} & \textbf{15{,}980} & \textbf{0.935} & \textbf{0.737} & \textbf{0.824} & \textbf{22}  \\
0.95 & 17{,}010 & 0.807 & 0.759 & 0.782 & 46  \\
0.92 & 18{,}301 & 0.495 & 0.785 & 0.607 & 123 \\
0.90 & 19{,}064 & 0.299 & 0.797 & 0.435 & 212 \\
0.88 & 19{,}744 & 0.228 & 0.806 & 0.355 & 261 \\
0.85 & 20{,}728 & 0.100 & 0.816 & 0.178 & 509 \\
0.80 & 22{,}405 & 0.048 & 0.832 & 0.091 & 762 \\
\bottomrule
\end{tabular}
\end{table}

% ============================================================================
\section{Results}
\label{sec:results}

\subsection{Quantitative}
Table~\ref{tab:main} reports our main numbers on the 30k held-out clustering
slice, plus the official Phase~1 leaderboard score.

\begin{table}[t]
\centering
\caption{Main results on the held-out 30k clustering slice (pairwise $F_1$),
and Phase~1 leaderboard score. ``Pair $F_1$'' refers to the binary pair
classifier; ``Cluster $F_1$'' is the Phase~2 metric.}
\label{tab:main}
\begin{tabular}{lccc}
\toprule
Method & Pair $F_1$ & Cluster $F_1$ & Phase-1 LB \\
\midrule
B1: Singleton            & --      & 0.000  & --     \\
B2: Exact-title          & --      & TBD    & --     \\
B3: Raw CLIP cosine      & --      & TBD    & --     \\
\textbf{Ours (baseline)} & \textbf{0.976} & \textbf{0.824} & \textbf{TBD} \\
\bottomrule
\end{tabular}
\end{table}

\subsection{Strengths}
The system handles \emph{cross-geo same-product matching} well: the
multilingual text encoder lets ``Sneakers Adidas Stan Smith'' (en) and
``Tenisky Adidas Stan Smith'' (cz) collapse to high cosine, and the price
feature (per-geo standardised) prevents trivial currency-induced false
negatives. Inference is cheap: $\approx 200$k items cluster in a few
minutes on an Apple~M4.

\subsection{Weaknesses}
The dominant failure mode is \emph{over-merging} on visually similar but
distinct fashion items---two different white sneakers, two different black
T-shirts. Both vanilla CLIP image embeddings and the raw cosine feature
saturate near $1.0$ for these confounders, and the Siamese head learns to
trust the saturated signal because the hard-negative mining cutoff
($s \ge 0.9$) does not penalise it. We confirmed this by sweeping
$\theta^\ast$: pair-$F_1$ stays $> 0.97$ across the sweep while clustering
$F_1$ collapses if $\theta^\ast$ is set below $\approx 0.95$, because a few
super-confident false-positive edges chain disjoint products into a single
giant connected component.

\subsection{Ablation}
Removing the raw $\langle \mathbf{e}^{\text{img}},
\mathbf{e}^{\text{img}}\rangle$ and $\langle \mathbf{e}^{\text{txt}},
\mathbf{e}^{\text{txt}}\rangle$ inputs to the classifier hurts pair-$F_1$
by $\approx 0.04$ in our preliminary runs but produces a noticeably more
spread-out logit distribution. We did not have time to retune
$\theta^\ast$ for the ablation and so do not report a clustering number.

% ============================================================================
\section{Conclusion}
\label{sec:conclusion}

A frozen-encoder Siamese architecture with raw modality-similarity features
and embedding-mined pairs is a surprisingly strong baseline for cross-geo
product matching. With only a $\approx 500$k-parameter head trained for ten
epochs on a laptop, we approach the strongest leaderboard scores reported
for this task, while remaining easy to debug and serve.

\textbf{Limitations and future work.}
(1) The hard-negative cutoff $s \ge 0.9$ is calibrated to the
\emph{training} similarity distribution, not the inference
KNN-candidate distribution; widening the band to $0.6 \le s < 1.0$ should
narrow the train/inference gap.
(2) Replacing vanilla CLIP with a fashion-domain encoder (FashionCLIP)
showed promising same/different-label cosine separation but did not
translate into clustering improvements within the current architecture, and
a metric-learning re-formulation (triplet/InfoNCE on the learned
embedding) is the natural next step.

% ============================================================================
\bibliographystyle{IEEEtran}
\bibliography{refs}

\end{document}
